\section{Finite-action covariant phase space with an inner boundary}\label{sec:finite-action-cps}

\subsection{Geometry, action, and orientations}\label{subsec:geometry-action}

Let $M_{\epsilon,R}$ be a regulated spacetime region bounded by initial and final Cauchy surfaces, an outer timelike cylinder $\Gamma_R$, and an inner timelike brick wall $\mathcal H_\epsilon$:

\begin{align}
\partial M_{\epsilon,R}
=\Sigma_f-\Sigma_i+\Gamma_R+\mathcal H_\epsilon .
\end{align}

Every timelike normal points out of the retained region. Thus the outer normal on $\Gamma_R$ points toward increasing radius, while the inner normal on $\mathcal H_\epsilon$ points into the excised collar. If the two cuts of a partial Cauchy surface are compared using the same reference tangent orientation, Stokes' theorem reads

\begin{align}
\partial\Sigma=C_R\sqcup(-C_\epsilon).
\end{align}

This is the origin of every outer-minus-inner sign below.

The bulk action is

\begin{align}
S_{\rm bulk}[g]
=\frac1{2\kappa_{\rm p}^{2}}
\int_{M_{\epsilon,R}}\!\mathrm d^3x\sqrt{-g}\,(R+2).
\end{align}

On a timelike component $\Gamma_a$ with outward unit normal $n_a$ and induced metric $\gamma^{(a)}_{ij}$, we use

\begin{align}
S_a
=\frac1{\kappa_{\rm p}^{2}}
\int_{\Gamma_a}\!\mathrm d^2x\sqrt{-\gamma^{(a)}}\,K_a
+S_{{\rm ct},a} .
\end{align}

The AdS counterterm $S_{{\rm ct},R}=-(1/\kappa_{\rm p}^{2})\int_{\Gamma_R}\sqrt{-\gamma}$ belongs to the asymptotic cylinder. It is not copied to the inner wall. For a fully Dirichlet piecewise-smooth region, non-orthogonal intersections also carry the Hayward action

\begin{align}
S_J=\frac{\sigma_J}{\kappa_{\rm p}^{2}}
\int_J\sqrt q\,\eta,
\end{align}

where $\eta$ is the relative boost angle and $\sigma_J$ is fixed by the induced joint orientation \cite{Hayward:1993my,Lehner:2016vdi}. The baseline CPS time-evolution problem and the fully Dirichlet problem must be kept distinct: in the former, the Cauchy surfaces are endpoints rather than boundary components with their own GHY terms.

\subsection{Componentwise first variation}\label{subsec:first-variation}

Write the Einstein--Hilbert potential as the two-form dual of

\begin{align}
\vartheta^\mu[g;\delta g]
=\frac1{2\kappa_{\rm p}^{2}}
\left(\nabla_\nu\delta g^{\mu\nu}-\nabla^\mu\delta g\right).
\end{align}

On each timelike component, the variation of the GHY term cancels the uncontrolled normal derivative of $\delta g$. The remaining local identity has the form

\begin{align}
\left.(\Theta_{\rm bulk}+\delta\ell_a)\right|_{\Gamma_a}
=\mathrm d C_a+\mathcal B_a,
\qquad
\mathcal B_a=-\frac12\Pi_a^{ij}\delta\gamma^{(a)}_{ij}, \label{eq:boundary-descent}
\end{align}

where

\begin{align}
\Pi_a^{ij}
=\frac{\sqrt{-\gamma^{(a)}}}{\kappa_{\rm p}^{2}}
\left(K_a^{ij}-K_a\gamma_{(a)}^{ij}
+\text{counterterm contribution}\right)
\end{align}

is the Brown--York momentum density. In the convention inherited from the asymptotic calculation, the endpoint one-form is

\begin{align}
C_a[\delta g]
=\int_{C_a}\sqrt q\,\tau_\mu C_a^\mu[\delta g],
\qquad
C_a^\mu[\delta g]
=-\frac1{2\kappa_{\rm p}^{2}}
\gamma_a^{\mu\nu}n_a^\rho\delta g_{\nu\rho}.
\end{align}

Equation (\ref{eq:boundary-descent}) separates three notions that are sometimes conflated. A Dirichlet boundary condition fixes $\delta\gamma_{ij}$. Gaussian radial gauge fixes components $n^\rho\delta g_{\rho\mu}$ and can set $C_a$ to zero. A zero-flux polarization requires the pullback of $\delta\mathcal B_a$ to vanish or be field-space exact. None of these conditions implies either of the other two.

If the chosen boundary conditions make the action differentiable, the integrated Cauchy-surface potential is

\begin{align}
\theta_\Sigma
=\int_\Sigma\Theta_{\rm bulk}
-C_R-C_\epsilon, \label{eq:integrated-potential}
\end{align}

and the pre-symplectic form is

\begin{align}
\Omega_\Sigma
=\delta\theta_\Sigma
=\int_\Sigma\omega_{\rm bulk}
-\delta C_R-\delta C_\epsilon . \label{eq:integrated-symplectic}
\end{align}

The signs in (\ref{eq:integrated-potential}) use the outward normal on each timelike boundary. If the inner expression is rewritten using a reference normal pointing toward increasing radius, its displayed sign reverses.

The Harlow--Wu one-form $C_a$ is not the Hayward scalar $S_J$. The former appears in the field-space descent (\ref{eq:boundary-descent}); the latter is part of a fully Dirichlet action and changes the endpoint potential by its variation. Their relation must be derived for a specified corner geometry rather than imposed by notation.

\subsection{Diffeomorphism descent and the Hamiltonian identity}\label{subsec:diffeomorphism-descent}

Let $\xi$ be a field-independent exact Killing field of the background that preserves the regulated boundaries and their declared sources. Acting on the complete action before imposing equations of motion gives

\begin{align}
X_\xi\mathbin{\cdot}\delta
\left(S_{\rm bulk}+S_R+S_\epsilon\right)
=\left.\alpha_\xi\right|_{\Sigma_f}
-\left.\alpha_\xi\right|_{\Sigma_i},
\end{align}

with

\begin{align}
\alpha_\xi
=\alpha_{\xi,{\rm bulk}}-\mu_{\xi,R}-\mu_{\xi,\epsilon}.
\end{align}

The finite-action generator is therefore

\begin{align}
H_\xi=X_\xi\mathbin{\cdot}\theta_\Sigma-\alpha_\xi,
\end{align}

and field-space Cartan calculus yields

\begin{align}
\delta H_\xi
=\Omega_\Sigma(\delta g,\mathcal L_\xi g). \label{eq:hamiltonian-identity}
\end{align}

Reducing the bulk current by Stokes' theorem and combining each cut with its own boundary descent gives

\begin{align}
\Omega_\Sigma(\delta g,\mathcal L_\xi g)
=\delta H_\xi^\infty-\delta H_\xi^\gamma
+\mathcal C_\xi[\delta g]
+\mathcal F_\xi[\delta g], \label{eq:outer-minus-inner}
\end{align}

where $\mathcal C_\xi$ denotes the bulk constraints and $\mathcal F_\xi$ the residual wall, joint, and source flux. Equation (\ref{eq:outer-minus-inner}) is an off-shell organization: the constraints and flux must not be discarded until their respective equations and boundary conditions have been imposed.

For the asymptotic component, the same calculation as in the parent finite-action construction combines the bulk surface form with $\mu_{\xi,R}-X_\xi\cdot C_R$ into the background-subtracted Brown--York charge \cite{Brown:1992br,Gao:2026graviton}. For the inner component, the corresponding expression is the finite-cutoff horizon charge, supplemented by any nonvanishing joint or embedding term required by the chosen variational problem.

\subsection{Perturbative order}\label{subsec:perturbative-order}

The first-order equation and the second-order equation are kept separate:

\begin{align}
\mathcal E^{(1)}[h]&=0,\\
\mathcal E^{(1)}[k]+\mathcal E^{(2)}[h,h]&=0.
\end{align}

The finite pure-$h$ current retains all three pieces

\begin{align}
H_{\xi,h}^{\mu}
=\xi_\nu T_h^{\mu\nu}
+\nabla_\nu S_{\xi,h}^{\mu\nu}
+\mathcal R_{\xi,h}^{\mu},
\qquad
T_h^{\mu\nu}=-\mathcal E^{(2),\mu\nu}[h,h].
\end{align}

The $k$-linear surface term and the $h^2$ term are combined only after using the second-order field equation. This bookkeeping is essential because the second derivative of a Hamiltonian along

\begin{align}
g(\lambda)=G+\lambda\kappa_{\rm p}h
+\lambda^2\kappa_{\rm p}^{2}k+O(\lambda^3)
\end{align}

is twice its $\lambda^2$ Taylor coefficient.
